\documentclass{article}
\usepackage{amsmath, amssymb, amsthm}
\usepackage{hyperref}
\usepackage{geometry}
\geometry{margin=1in}

\title{Erd\H{o}s Problem \#857}
\date{Last updated: 2025-08-31}
\author{Source: erdosproblems.com}

\begin{document}
\maketitle

\noindent\textbf{Status:} OPEN \quad \textbf{No prize}

\noindent\textbf{Tags:} combinatorics

\medskip
\noindent\textbf{Note:} weak sunflower problem

\medskip

\noindent\textbf{Problem Statement:}

Let $m=m(n,k)$ be minimal such that in any collection of sets $A_1,\ldots,A_m\subseteq \{1,\ldots,n\}$ there must exist a sunflower of size $k$ - that is, some collection of $k$ of the $A_i$ which pairwise have the same intersection. Estimate $m(n,k)$, or even better, give an asymptotic formula.




Related to [536] and [856] . In \cite{Er70} Erd\H{o}s asks this in the equivalent formulation with intersection replaced by union. This is sometimes known as the weak sunflower problem (see [20] for the strong sunflower problem). When $k=3$ this is strongly connected to the cap set problem (finding the maximal size of subsets of $\mathbb{F}_3^n$ with no three-term arithmetic progressions), as observed by Alon, Shpilka, and Umans \cite{ASU13}). Naslund and Sawin \cite{NaSa17} have proved that\[m(n,3) \leq (3/2^{2/3})^{(1+o(1))n}.\]




References

[ASU13] Alon, Noga and Shpilka, Amir and Umans, Christopher, On sunflowers and matrix multiplication . Comput. Complexity (2013), 219--243.

[Er70] Erd\H{o}s, Paul, Some extremal problems in combinatorial number theory . Mathematical Essays Dedicated to A. J. Macintyre (1970), 123-133.

[NaSa17] Naslund, Eric and Sawin, Will, Upper bounds for sunflower-free sets . Forum Math. Sigma (2017), Paper No. e15, 10.

\noindent\textbf{Related OEIS sequences:} possible


\bigskip
\noindent\small{Source: \url{https://www.erdosproblems.com/857}}
\end{document}
